\documentclass[10pt,a4paper]{article}
% ========================== DO NOT MODIFY =================================== %
\usepackage[a4paper,left=2.0cm,top=2.0cm,right=2.0cm,bottom=2.0cm]{geometry}
\usepackage{amssymb,amsmath} % If amsmath is required
\usepackage[implicit=false]{hyperref}
\usepackage[utf8]{inputenc}
\usepackage[T1]{fontenc}
\usepackage{times}
\usepackage{graphicx}
\usepackage{titling}
\usepackage{authblk}
\usepackage{multicol}
\usepackage{float}
\usepackage{enumitem}
\usepackage{xcolor}
\usepackage{hyperref}
\hypersetup{
    colorlinks=true,
    linkcolor=blue,
    filecolor=blue,      
    urlcolor=blue,
}
    
\usepackage{fancyhdr}
\fancyhf{}
\setlength{\headheight}{3em}
\pagestyle{fancy}
\renewcommand{\headrulewidth}{0pt} % remove lines as well
\fancyhead[L]{\footnotesize{23\textsuperscript{st} International Conference on Biomedical Applications of Electrical Impedance Tomography (EIT 2023), Aachen, Germany}}
\fancypagestyle{firstpage}{%
  \lhead{\footnotesize{23\textsuperscript{st} International Conference on Biomedical Applications of Electrical Impedance Tomography (EIT 2023), Aachen, Germany}}
}

\usepackage[small,compact]{titlesec}
% Fix titlesec for the 2.10.1 version in ubuntu 16.04
%  https://bugs.launchpad.net/ubuntu/+source/texlive-extra/+bug/1574052
\makeatletter
\@ifpackagelater{titlesec}{2016/03/21}{%
 % Package titlesec is on version 2.10.2 or higher, nothing to do %
}{%
 % Check if package titlesec is on version 2.10.1 %
 \@ifpackagelater{titlesec}{2016/03/15}{%
  % Package titlesec on version 2.10.1, patch accordingly %
  \usepackage{etoolbox}%
  \makeatletter
  \patchcmd{\ttlh@hang}{\parindent\z@}{\parindent\z@\leavevmode}{}{}%
  \patchcmd{\ttlh@hang}{\noindent}{}{}{}%
  \makeatother
 }{%
  % Package titlsecon is on version 2.10.0 or lower, nothing to do %
 }%
}
\makeatother
\titlespacing*{\section}{0pt}{5pt}{5pt}  % If you need different spacing

\usepackage[numbers,sort&compress]{natbib}
\usepackage[small,bf,tableposition=top,figureposition=bottom,skip=2pt]{caption}


% limit space between floats and text
\setlength{\textfloatsep}{10pt plus 1.0pt minus 2.0pt}
% shrink the list environments
\setlist{noitemsep}
\setlist{nolistsep}

% Shrink the bibliography
  \let\oldthebibliography=\thebibliography
  \let\endoldthebibliography=\endthebibliography
  \renewenvironment{thebibliography}[1]{%
    \begin{oldthebibliography}{#1}%
      \setlength{\parskip}{0ex}%
      \setlength{\itemsep}{0ex}%
  }%
  {%
    \end{oldthebibliography}%
  }

% Format of the front matter
\renewcommand{\Affilfont}{\small}
\setlength{\affilsep}{1ex}
\renewcommand\Authands{, }
\setlength{\droptitle}{-2cm}
% ============================================================================ %




%------------------------------- FRONT MATTER -------------------------------- %
\title{HDF5-based data format for EIT data%
\vspace{-0.5ex}} %remove vertical space
\author[1]{Lucas Possner}
\author[1]{Martin Bulst}
\author[2]{Andy Adler}
\affil[1]{Sciospec Scientific Instruments GmbH, Leipzig, Germany}
\affil[2]{Carleton University, Ottawa, Canada}
\date{}
%----------------------------------------------------------------------------- %
%\protect\url


\begin{document}
\maketitle
\thispagestyle{firstpage}
\vspace{-1.5cm}
%\thispagestyle{empty}

\begin{multicols}{2}

\noindent {\bf Abstract:}
Many devices, both research and commercial, have been
made for EIT and bioimpedance applications. We propose
a new, flexible data format which can allow data
interchange, based on the HDF5 data standard.

\section{Introduction}
EIT and bioimpedance applications are growing. For 
research applications, there is an increasing
complexity of configurations in which complex
configurations and multi data streams are used.
We have previously proposed an openEIT data format\cite{Jones2014Open},
which received little traction because of the difficulties
of using XML as a data container.
%
For a new data format, our requirements are:
\begin{itemize}[itemsep=0ex, leftmargin=5mm,topsep=2mm]
\item Simplicity:
  writing data in $<50$ code lines. 
\item Self-documenting: 
  core data element is the ``frame''.
\item Streaming: suitable for large
  and online data sets.
\item Extensibility:
  easy to add additional data and elements.
\end{itemize}
Here, we propose a data format based on HDF5 \cite{HDF5Group},
a hierarchical format designed for numerical data.


%HDF5 stores data in row major order:

% HDF5 uses C storage conventions, assuming that the last listed
% dimension is the fastest-changing dimension and the first-listed
% dimension is the slowest changing

\section{Format}

The format (Version 2023.4) is illustrated in Fig\,\ref{fig:system};
updated versions will be documented at \cite{Adler2023HDF5}.
Data sets are stored as frames of data, where each is a sequence
of measurements,
which are repeated over time. For complex data, 
{\tt Meas.V.Real}
and 
{\tt Meas.V.Imag}
data sets are stored.
 Multiple data sets can be stored in one file,
and each measurement configuration is documented in a 
{\tt /data/\emph{\{Data set name\}}/protocol} field.
Electrodes can be numbered or named. In the
protocol, the field
 {\tt Stim.I.01(A)} indicates the excitation current into electrode {\tt 01} in Amps;
 the $4^{\rm th}$ element shows $10$\,mA=$0.01$\,A is injected
during the 
     $4^{\rm th}$ measurement in each frame.
Excitation (stimulation) information is stored into fields that start with
{\tt Stim} and specify the kind of excitation after a dot.
Similary, measurement information are stored into fields that start with {\tt Meas}.
Measurements normally correspond to stimulations, but this
is not required e.g.\ for biopotential measurements, where
{\tt Meas.V.freq(Hz)}=0.
Measurement data are stored as a matrix
 $N_{\rm frames}\times L_{\rm frame}$,
with an extension indicating the type
({\tt .Abs}, {\tt .Real}, {\tt .Imag}).
Times (if known) can be stored as shown.
Optional / vendor-specific information
 (e.g.\ {\tt /instrument}, or
       {\tt /data/initialize})
can be stored.
Additional recordings with their own protocols are stored with different 
{\tt /data/\emph{\{Data set name\}}}.

\section{Discussion}

We describe an HDF5-based data format for EIT data.
It documented with code examples at \cite{Adler2023HDF5}.
The data format is supported in EIDORS \cite{Adler2006EIDORS}
and will be part of the next software release.


\footnotesize
% Instead of using thebibliography, you can also type the references
% using bibtex, and uncomment the following lines
%\bibliographystyle{compact}
%\bibliography{eit2018}
\begin{thebibliography}{}
\bibitem{Jones2014Open}
C Jones \emph{et al},
%Open Electrical Impedance Tomography (OEIT) File Format
\href{https://cudo.carleton.ca/eit2014/3311}{p.20}, Conf EIT 2014
%http://www.sce.carleton.ca/faculty/adler/publications/2014/jones-EIT2014-OpenEIT.pdf

\bibitem{Adler2006EIDORS}
A Adler, WRB Lionheart, 
%EIDORS
Physiol Meas 27:S25--42, 2006

\bibitem{Adler2023HDF5}
A Adler \emph{et al},
\href{https://sourceforge.net/p/eidors3d/code/HEAD/tree/trunk/dev/HDF5\_tools}
                      {sf.net/p/eidors3d/code/HEAD/tree/trunk/dev/HDF5\_tools/}
``HDF5 tools", 2023

\bibitem{HDF5Group}
The HDF Group,
%    title = "{Hierarchical data format version 5}",
1997--2023,
\href{http://www.hdfgroup.org/HDF5}
                {hdfgroup.org/HDF5}

\end{thebibliography}
\end{multicols}

\vspace{-5mm}

\begin{figure*}[h]
\centering
% Use the following line with your images
\includegraphics[page=13,width=.92\textwidth]{figures.pdf}
\vspace{-1mm}
\caption{
Structure of the proposed HDF5 file (Version 2023.4).
All data sets must contain
a {\tt protocol} defining the configuration for each measurement
in a frame. Frame data is a single file (e.g.\ {\tt Meas.V.abs}).
Times for measurements or frames are provided if available.
Items in {\em italics} have example names (can be named differently).
}
\label{fig:system}
\end{figure*}

\end{document}
